
\documentclass{article}
\usepackage{colortbl}
\usepackage{makecell}
\usepackage{multirow}
\usepackage{supertabular}

\begin{document}

\newcounter{utterance}

\centering \large Interaction Transcript for game `referencegame', experiment `letter\_grids', episode 11 with reflexion{-}Qwen3.5{-}9B.
\vspace{24pt}

{ \footnotesize  \setcounter{utterance}{1}
\setlength{\tabcolsep}{0pt}
\begin{supertabular}{c@{$\;$}|p{.15\linewidth}@{}p{.15\linewidth}p{.15\linewidth}p{.15\linewidth}p{.15\linewidth}p{.15\linewidth}}
    \# & $\;$A & \multicolumn{4}{c}{Game Master} & $\;\:$B\\
    \hline

    \theutterance \stepcounter{utterance}  
    & & \multicolumn{4}{p{0.6\linewidth}}{
        \cellcolor[rgb]{0.9,0.9,0.9}{
            \makecell[{{p{\linewidth}}}]{
                \texttt{\tiny{[P1$\langle$GM]}}
                \texttt{You have attempted this same task before and did not succeed. These are the plans you wrote after those attempts.} \\
\texttt{Trial \#1: New plan: My previous description failed because it incorrectly identified the "bottom{-}left column" as completely filled (it is only filled in the first three rows) and mischaracterized the empty area as a "2x4 area" (the empty region is actually a 3x4 block spanning rows 2{-}4 and columns 2{-}5). To succeed, I must strictly adhere to the visual evidence: describe the top row as fully filled, the leftmost column as filled only for the first three rows, and the third row as fully filled. I will explicitly state that the remaining cells (rows 2{-}4, columns 2{-}5) are empty to ensure the description uniquely distinguishes the target from the distractors, which have different patterns of filled rows and columns.} \\
\\ 
\\ 
\texttt{You are given three grids, where each of them is 5 by 5 in size.} \\
\texttt{Grids have empty cells marked with "▢" and filled cells marked with "X".} \\
\texttt{Your task is to generate a referring expression that best describes the target grid while distinguishing it from the two other distractor grids.} \\
\texttt{The first grid is the target grid, and the following two grids are the distractors.} \\
\\ 
\texttt{Target grid:} \\
\\ 
\texttt{X X X X X} \\
\texttt{□ □ □ □ □} \\
\texttt{X X X X □} \\
\texttt{X □ □ □ □} \\
\texttt{X □ □ □ □} \\
\\ 
\texttt{Distractor grid 1:} \\
\\ 
\texttt{X X X X X} \\
\texttt{X □ □ □ □} \\
\texttt{X X X X X} \\
\texttt{X □ □ □ □} \\
\texttt{X X X X X} \\
\\ 
\texttt{Distractor grid 2:} \\
\\ 
\texttt{X X X X X} \\
\texttt{X □ □ □ □} \\
\texttt{X □ □ □ □} \\
\texttt{X □ □ □ □} \\
\texttt{X X X X X} \\
\\ 
\texttt{Instruction: Describe the target grid.} \\
\texttt{Generate the referring expression starting with the tag "Expression: " for the given target grid. Omit any other text.} \\
            }
        }
    }
    & \\ \\

    \theutterance \stepcounter{utterance}  
    & \multicolumn{4}{p{0.6\linewidth}}{
        \cellcolor[rgb]{0.9,0.9,0.9}{
            \makecell[{{p{\linewidth}}}]{
                \texttt{\tiny{[P1$\rangle$GM]}}
                \texttt{Expression: The grid has a completely filled top row and a completely filled third row; the leftmost column is filled only in the first three rows (rows 1, 2, and 3), while all other cells in the leftmost column (rows 4 and 5) are empty, and the remaining cells in rows 2, 3, and 4 (columns 2 through 5) are empty.} \\
            }
        }
    }
    & & \\ \\

    \theutterance \stepcounter{utterance}  
    & & & \multicolumn{2}{p{0.3\linewidth}}{
        \cellcolor[rgb]{0.9,0.9,0.9}{
            \makecell[{{p{\linewidth}}}]{
                \texttt{\tiny{[GM$|$GM]}}
                \texttt{Expression: The grid has a completely filled top row and a completely filled third row; the leftmost column is filled only in the first three rows (rows 1, 2, and 3), while all other cells in the leftmost column (rows 4 and 5) are empty, and the remaining cells in rows 2, 3, and 4 (columns 2 through 5) are empty.} \\
            }
        }
    }
    & & \\ \\

    \theutterance \stepcounter{utterance}  
    & & \multicolumn{4}{p{0.6\linewidth}}{
        \cellcolor[rgb]{0.9,0.9,0.9}{
            \makecell[{{p{\linewidth}}}]{
                \texttt{\tiny{[GM$\rangle$P2]}}
                \texttt{You have attempted this same task before and did not succeed. These are the plans you wrote after those attempts.} \\
\texttt{Trial \#1: New plan: I incorrectly identified the "Second" grid because I focused on the filled top and bottom rows but failed to verify the specific "completely filled third row" requirement. The Second grid has a gap in the third row (the middle cell is empty), whereas the Third grid has a completely filled third row (except for the last cell which is part of the empty bottom{-}right area, but the expression specifies the third row is filled, which matches the Third grid's pattern of X X X X □ where the first four are filled? Wait, let's re{-}read the expression carefully: "completely filled third row". Looking at the Third grid: Row 3 is "X X X X □". That is not completely filled. Let's re{-}examine the grids.} \\
\texttt{Grid 1: Row 3 is "X □ □ □ □". Not filled.} \\
\texttt{Grid 2: Row 3 is "X X X X X". Completely filled.} \\
\texttt{Grid 3: Row 3 is "X X X X □". Not completely filled.} \\
\texttt{Wait, the expression says "leaving the bottom{-}right 2x4 area empty".} \\
\texttt{Let's re{-}evaluate Grid 2 against the full expression:} \\
\texttt{1. Completely filled top row: Yes (X X X X X).} \\
\texttt{2. Completely filled bottom{-}left column: Yes (X in col 1 for all rows).} \\
\texttt{3. Completely filled third row: Yes (X X X X X).} \\
\texttt{4. Bottom{-}right 2x4 area empty: Rows 4 and 5, Cols 2{-}5.} \\
\texttt{   Grid 2 Row 4: X □ □ □ □ (Cols 2{-}5 are empty).} \\
\texttt{   Grid 2 Row 5: X X X X X (Cols 2{-}5 are NOT empty).} \\
\texttt{So Grid 2 fails the "bottom{-}right 2x4 area empty" condition because Row 5 is full.} \\
\\ 
\texttt{Let's re{-}evaluate Grid 3:} \\
\texttt{1. Top row: X X X X X (Yes).} \\
\texttt{2. Bottom{-}left column: Row 1(X), Row 2(□), Row 3(X), Row 4(X), Row 5(X). No, Row 2 is empty. So Grid 3 fails the "completely filled bottom{-}left column" condition.} \\
\\ 
\texttt{Let's re{-}evaluate Grid 1:} \\
\texttt{1. Top row: X X X X X (Yes).} \\
\texttt{2. Bottom{-}left column: All X (Yes).} \\
\texttt{3. Third row: X □ □ □ □ (No, not completely filled).} \\
\\ 
\texttt{There seems to be a contradiction in my analysis or the grids provided in the prompt text vs the description. Let's look closer at the prompt's Grid 2 again.} \\
\texttt{Grid 2:} \\
\texttt{Row 1: X X X X X} \\
\texttt{Row 2: X □ □ □ □} \\
\texttt{Row 3: X X} \\
\\ 
\\ 
\texttt{You are given three grids, where each of them is 5 by 5 in size.} \\
\texttt{Grids have empty cells marked with "▢" and filled cells marked with "X".} \\
\texttt{You are also given a referring expression that describes one of the given grids.} \\
\texttt{Your task is to select the grid that matches the given referring expression.} \\
\texttt{Generate only the number (in text) of the grid that matches the given expression by selecting first, second, or third.} \\
\\ 
\texttt{First:} \\
\\ 
\texttt{X X X X X} \\
\texttt{X □ □ □ □} \\
\texttt{X □ □ □ □} \\
\texttt{X □ □ □ □} \\
\texttt{X X X X X} \\
\\ 
\texttt{Second:} \\
\\ 
\texttt{X X X X X} \\
\texttt{X □ □ □ □} \\
\texttt{X X X X X} \\
\texttt{X □ □ □ □} \\
\texttt{X X X X X} \\
\\ 
\texttt{Third:} \\
\\ 
\texttt{X X X X X} \\
\texttt{□ □ □ □ □} \\
\texttt{X X X X □} \\
\texttt{X □ □ □ □} \\
\texttt{X □ □ □ □} \\
\\ 
\texttt{Expression: The grid has a completely filled top row and a completely filled third row; the leftmost column is filled only in the first three rows (rows 1, 2, and 3), while all other cells in the leftmost column (rows 4 and 5) are empty, and the remaining cells in rows 2, 3, and 4 (columns 2 through 5) are empty.} \\
\texttt{Question: Which grid does the expression refer to?} \\
\texttt{Start with the tag "Answer: ", followed by your selection. Omit any other text.} \\
            }
        }
    }
    & \\ \\

    \theutterance \stepcounter{utterance}  
    & & & \multicolumn{4}{p{0.6\linewidth}}{
        \cellcolor[rgb]{0.9,0.9,0.9}{
            \makecell[{{p{\linewidth}}}]{
                \texttt{\tiny{[GM$\langle$P2]}}
                \texttt{Answer: First} \\
            }
        }
    }
     \\ \\

    \theutterance \stepcounter{utterance}  
    & & & \multicolumn{2}{p{0.3\linewidth}}{
        \cellcolor[rgb]{0.9,0.9,0.9}{
            \makecell[{{p{\linewidth}}}]{
                \texttt{\tiny{[GM$|$GM]}}
                \texttt{first} \\
            }
        }
    }
    & & \\ \\

    \theutterance \stepcounter{utterance}  
    & \multicolumn{4}{p{0.6\linewidth}}{
        \cellcolor[rgb]{0.95,0.95,0.95}{
            \makecell[{{p{\linewidth}}}]{
                \texttt{\tiny{[P2$\rightarrow$P1]}}
                \texttt{Post{-}episode Reflexion (not sent to GM):} \\
\texttt{New plan: My previous description failed because it incorrectly stated that the leftmost column is filled only in the first three rows; in reality, the leftmost column is completely filled (rows 1{-}5). Additionally, I misidentified the empty region as spanning rows 2{-}4, when it actually spans rows 2{-}5. To succeed, I must strictly adhere to the visual evidence: describe the top row (row 1) and the third row (row 3) as completely filled, the leftmost column as completely filled, and explicitly state that the remaining cells (rows 2, 4, and 5, columns 2 through 5) are empty. This precise distinction will differentiate the target from Distractor 1 (which has a filled second row) and Distractor 2 (which has a filled fifth row).} \\
            }
        }
    }
    & & \\ \\

    \theutterance \stepcounter{utterance}  
    & & & \multicolumn{4}{p{0.6\linewidth}}{
        \cellcolor[rgb]{0.95,0.95,0.95}{
            \makecell[{{p{\linewidth}}}]{
                \texttt{\tiny{[P1$\rightarrow$P2]}}
                \texttt{Post{-}episode Reflexion (not sent to GM):} \\
\texttt{New plan: I failed because I relied on a partial visual match (filled top row) without rigorously verifying every specific constraint in the expression, particularly the "completely filled third row" and the "leftmost column filled only in the first three rows" conditions. My previous analysis was confused and contradictory. Next time, I will systematically check each grid against every clause of the expression in order: 1) Verify top row is full. 2) Verify third row is full. 3) Verify the leftmost column has Xs only in rows 1{-}3 and empty in rows 4{-}5. 4) Verify the remaining specified empty areas. Only after confirming all conditions are met will I select the grid. Based on this strict check: Grid 1 fails (Row 3 not full). Grid 2 fails (Row 5 is full, violating the empty bottom{-}right area and left column constraints). Grid 3 fails (Row 2 left cell is empty, violating the left column constraint). Wait, re{-}reading the expression: "leftmost column is filled only in the first three rows". Grid 3 Row 2 Col 1 is "□", which matches "empty". Grid 3 Row 1 Col 1 is "X", Row 3 Col 1 is "X", Row 4 Col 1 is "X", Row 5 Col 1 is "X". This violates "only in the first three rows". Let's re{-}read the expression carefully: "leftmost column is filled only in the first three rows (rows 1, 2, and 3), while all other cells in the leftmost column (rows 4 and 5) are empty".} \\
\texttt{Grid 1: Left col is X,X,X,X,X. Fails (rows 4,5 not empty).} \\
\texttt{Grid 2: Left col is X,X,X,X,X. Fails (rows 4,5 not empty).} \\
\texttt{Grid 3: Left col is X,□,X,X,X. Fails (Row 2 is empty, Row 4,5 are X).} \\
\texttt{There seems to be no grid matching the description perfectly as written, or I am misinterpreting the grid text. Let's look at Grid 3 again.} \\
\texttt{Grid 3:} \\
\texttt{Row 1: X X X X X} \\
\texttt{Row 2: □ □ □ □ □ (Left is empty)} \\
\texttt{Row 3: X X X X □} \\
\texttt{Row 4: X □ □ □ □} \\
\texttt{Row 5: X □ □ □ □} \\
\texttt{Expression check for Grid 3:} \\
\texttt{{-} Top row full? Yes.} \\
\texttt{{-} Third row full? No (last cell is □).} \\
\texttt{{-} Left col filled only in rows 1{-}3? No (Row 2 is empty).} \\
\texttt{{-} Bottom{-}right 2x4 empty? Rows 4{-}5, Cols 2{-}5 are empty} \\
            }
        }
    }
     \\ \\

\end{supertabular}
}

\end{document}
